\clearpage

\section {Introduction }
\label{sec:intro}

In this note we present the results of a 3-flavour $\nu_{\mu}$-disappearance analysis on the combined T2K Run1+2
dataset. This work builds on the effective 2-flavour $\nu_{\mu}$-disappearance analysis of the Run1+2 dataset
that was reported in Ref.\ \cite{OurRun1and2Analysis2F}.
Further technical details of the effective 2-flavour analysis may be found in the analysis of the
Run 1 dataset that was reported in Ref.\ \cite{OurRun1DisappAnalysis2F}.

As the measurement error on the atmospheric $\Delta m^{2}$ 
approaches the solar $\Delta m^{2}$ ($\sim$0.075$\times$10$^{-3}$ eV$^2$/c$^{4}$),
subleading effects of the solar $\Delta m^{2}$ on $\nu_{\mu}$-disappearance analyses should not be neglected. 
This was pointed out by E.Lisi in the XIV International Workshop on Neutrino Telescopes \cite{LisiNeuTel11},
and brought to the attention of the T2K $\nu_{\mu}$-disappearance contour committee by M. Mezzetto.
In response, the contour committee recommended moving to a 3-flavour analysis \cite{T2KNuMuContourCom}.

The method for predicting the nominal Super-K reconstructed energy spectrum of single $\mu$-like ring events 
is described in Ref.\ \cite{OurRun1and2Analysis2F} and it is unchanged.
This work considers the same systematic uncertainties as Ref.\ \cite{OurRun1and2Analysis2F}.
As in the effective 2-flavour analysis, the best-fit oscillation parameters were found using the likelihood-ratio method \cite{StatAnalysisCowan}, 
and the confidence regions were constructed using the Feldman-Cousins method \cite{FeldmanCousins}.
Systematic errors were included in the oscillation contours using the Cousins-Highland method (Bayesian integration) 
and by determining their effects on the Feldman-Cousins $\Delta\chi^{2}_{critical}$ values.

The oscillation probabilities in this work were calculated in a full 3-flavour framework including matter effects. 
The $\nu_{\mu}$-survival probability depends on sin$^{2}$2$\theta_{12}$, sin$^{2}$2$\theta_{13}$, $\Delta m^{2}_{12}$, $\delta_{CP}$ 
and the neutrino mass hierarchy as well as sin$^{2}$2$\theta_{23}$ and $|\Delta m^{2}_{32}|$. Throughout this note $\sin^{2}(2 \theta_{12})$, 
$\sin^{2}(2 \theta_{13})$, $\Delta m_{21}^{2}$, $\delta_{CP}$ will be collectively denoted the ``non-23'' oscillation parameters. The method of 
calculation of the 3-flavour oscillation probabilities is described in Sec.\ \ref{sec:nuosc3f}, and details of the validation of the code and our 
assessment of the its numerical accuracy are also presented. The predicted reconstructed single $\mu$-like ring spectra in 
Super-K for the 3-flavour analysis are shown in Sec. \ref{sec:pred}. Details of the 3-flavour fits of toy MC datasets are given in 
Sec. \ref{sec:fits}; in this analysis, sin$^{2}$2$\theta_{12}$, sin$^{2}$2$\theta_{13}$, $\Delta m^{2}_{21}$ and $\delta_{CP}$ were 
randomised when creating the toy MC datasets but were fixed when fitting them or the real data. This section also includes the results of the 3-flavour 
fit of the Run1+2 dataset for both GENIE and NEUT, and a comparison with the results of the 2-flavour analysis \cite{OurRun1and2Analysis2F}. 
The 3-flavour confidence regions are shown in Sec.\ \ref{sec:conf_regions}, and these are compared with those obtained 
in the 2-flavour analysis \cite{OurRun1and2Analysis2F} and with previous results from MINOS \cite{Minos2011} and Super-K \cite{OldDisappSuperK}. 
A summary is given in Sec.\ \ref{sec:summary}.  

